El informe documenta el desarrollo, calibraci\'on y evaluaci\'on de negocio de un
sistema recomendador de inversiones para FinPUC, una plataforma ficticia de
asesor\'ia financiera que genera portafolios personalizados para clientes del
mercado accionario estadounidense. El sistema atiende cinco perfiles de riesgo,
desde muy conservador hasta muy arriesgado, y persigue un \'unico objetivo:
maximizar el retorno esperado del cliente. La utilidad de la empresa se mide
\emph{a posteriori}, como consecuencia del saldo activo administrado, mediante
una comisi\'on mensual fija de 0,5\% sobre dicho saldo, sin comisiones por
transacci\'on ni por rebalanceo.

El universo de trabajo comprende 601 acciones, obtenidas tras filtrar 1\,406
tickers. Se implementaron tres capas metodol\'ogicas: optimizaci\'on
media-varianza de Markowitz como caso base, cuatro familias de \textit{views} de
Black-Litterman calibradas de forma independiente, y una simulaci\'on Monte Carlo
que modela el comportamiento del cliente. Cada \textit{view} se calibra bajo un
protocolo secuencial de tres horizontes (calibraci\'on, validaci\'on y test
limpio) para no contaminar el per\'iodo de evaluaci\'on. El comportamiento del
cliente se reduce a una \'unica probabilidad de abandono semanal; las
recomendaciones se emiten de forma semestral, mientras que el abandono se
eval\'ua semanalmente.

Entre las \textit{views} calibradas, momentum general entrega la mayor ganancia
econ\'omica ($+85.7\%$ en Score P4 y $+5.2\%$ en Sharpe frente a Markowitz base,
cifras promedio), aunque con mayor rotaci\'on (34.6\%) y concentraci\'on
sectorial (HHI 0.261), mientras que desempleo macro ofrece la mejor estabilidad
operacional (7.3\% de turnover, HHI 0.157). La metodolog\'ia recomendada es
Black-Litterman con la \textit{view} de momentum general para perfiles agresivos,
con desempleo macro como alternativa defensiva para perfiles conservadores.

\textbf{Palabras clave:} optimizaci\'on de portafolios, Markowitz,
Black-Litterman, calibraci\'on independiente, simulaci\'on Monte Carlo, sistema
recomendador, abandono semanal, comisi\'on sobre saldo administrado, validaci\'on
rodante, CVXPY.
